\documentclass[12pt]{article}
\usepackage[utf8]{inputenc}
\usepackage[margin=1in]{geometry}
\usepackage{amsmath, amssymb, amsthm}
\usepackage{xcolor}
\usepackage{tcolorbox}
\usepackage{enumitem}
\usepackage{fancyhdr}
\usepackage{titlesec}

% Define colored boxes
\newtcolorbox{strategybox}[1]{
	colback=blue!5!white,
	colframe=blue!75!black,
	title=#1,
	fonttitle=\bfseries
}

\newtcolorbox{tacticsbox}[1]{
	colback=green!5!white,
	colframe=green!75!black,
	title=#1,
	fonttitle=\bfseries
}

\newtcolorbox{toolsbox}[1]{
	colback=red!5!white,
	colframe=red!75!black,
	title=#1,
	fonttitle=\bfseries
}

\newtcolorbox{examplebox}[1]{
	colback=orange!5!white,
	colframe=orange!75!black,
	title=#1,
	fonttitle=\bfseries
}

\newcommand{\docversion}{2.0.1}

\title{\textbf{Strategies and Tactics Guide:\\The Art and Craft of Problem Solving by Paul Zeitz}}
\author{Comprehensive Educational Resource}
\date{\today \space \docversion}
\begin{document}
	
	\maketitle
	
	\tableofcontents
	\newpage
	
	\section{Introduction and Core Philosophy}
	
	Paul Zeitz's \textit{The Art and Craft of Problem Solving} presents a revolutionary three-tier hierarchical framework for mathematical problem solving. The book distinguishes between \textbf{exercises} (immediately solvable using known techniques) and \textbf{problems} (requiring investigation and resourcefulness).
	
	\begin{strategybox}{Core Pedagogical Philosophy}
		\textbf{Three-Level Problem Solving Hierarchy:}
		\begin{enumerate}
			\item \textbf{Strategy:} Mathematical and psychological ideas for starting and pursuing problems
			\item \textbf{Tactics:} Diverse mathematical methods that work in many different settings
			\item \textbf{Tools:} Narrowly focused techniques and "tricks" for specific situations
		\end{enumerate}
		
		\textbf{Key Principle:} "Anything that stimulates investigation is good" (p. 23)
	\end{strategybox}
	
	\section{Strategic Framework Analysis}
	
	\subsection{Psychological Strategies}
	
	\subsubsection{Mental Toughness and Creativity (Chapter 2.1)}
	
	\begin{strategybox}{Mental Toughness and Creativity}
		\textbf{Mental Toughness - Pólya's Mouse Principle:}
		\begin{itemize}
			\item Persistence through multiple failed attempts
			\item "Toughen up, loosen up, and practice" (p. 20)
			\item Avoid immediate declarations of impossibility
		\end{itemize}
		
		\textbf{Creativity Cultivation:}
		\begin{itemize}
			\item "Learn to shamelessly appropriate new ideas and make them your own" (p. 16)
			\item Develop mathematical peripheral vision
			\item Break or bend rules deliberately
			\item "Good, obedient boys and girls solve fewer problems than naughty and mischievous ones" (p. 19)
		\end{itemize}
	\end{strategybox}
	
	\subsubsection{Iterative Refinement Through Feedback Loops (Chapter 2.1)}
	
	\begin{strategybox}{Iterative Refinement}
		\textbf{Iterative Approach:}
		\begin{itemize}
			\item Start with rough estimates or simplified versions of the problem
			\item Refine solutions through successive approximations based on feedback
			\item View errors as opportunities to adjust and improve approaches
		\end{itemize}
		
		\textbf{Application:}
		\begin{itemize}
			\item Builds resilience by normalizing iterative failure and correction
			\item Encourages adaptability in problem-solving mindset
			\item Complements mental toughness by emphasizing persistence in refinement
		\end{itemize}
	\end{strategybox}
	
	\subsection{Investigation Strategies}
	
	\subsubsection{Startup Orientation (Chapter 2.2)}
	
	\begin{strategybox}{Startup Orientation}
		\textbf{Orientation Process:}
		\begin{enumerate}
			\item Read problem carefully, paying attention to details
			\item Classify as "to find" or "to prove" problem
			\item Identify hypothesis and conclusion clearly
			\item Brainstorm notation, methods, and intuitive solutions
		\end{enumerate}
	\end{strategybox}	

	\subsubsection{Get Your Hands Dirty (Chapter 2.2)}

	\subsubsection{Penultimate Step (Chapter 2.2)}

	\subsubsection{Wishful Thinking (Chapter 2.2)}

	\subsubsection{Make It Easier (Chapter 2.2)}

	\subsubsection{Conjecture via Small Cases }

	\subsubsection{Visualization and Diagramming for Insight (Chapter 2.2)}
	
	\begin{strategybox}{Visualization and Diagramming}
		\textbf{Visual Representation:}
		\begin{itemize}
			\item Create diagrams, graphs, or sketches to represent abstract concepts
			\item Use freehand drawings or software for geometry and graph theory problems
			\item Identify patterns or symmetries through visual aids
		\end{itemize}
		
		\textbf{Benefits:}
		\begin{itemize}
			\item Translates numerical or algebraic issues into visual forms
			\item Facilitates hypothesis formation and solution discovery
			\item Enhances understanding of complex relationships
		\end{itemize}
	\end{strategybox}
	
	\subsection{Argument Construction}
	
	\subsubsection{Direct Proof}

	\begin{strategybox}{Direct Proof}
		\textbf{Direct Proof:} Straightforward deductive reasoning
	\end{strategybox}

	\subsubsection{Proof by Contradiction}

	\begin{strategybox}{Proof by Contradiction}
		\textbf{Proof by Contradiction:} Assume negation and derive contradiction
	\end{strategybox}

	\subsubsection{Mathematical Induction}

	\begin{strategybox}{Mathematical Induction}
		\textbf{Mathematical Induction:} Base case + inductive step
	\end{strategybox}

	\subsubsection{Stylistic Conventions}

	\begin{strategybox}{Stylistic Conventions}
		\textbf{Stylistic Conventions:}
		\begin{itemize}
			\item Clear hypothesis and conclusion statements
			\item WLOG (without loss of generality) for similar cases
			\item TS/ISTS (to show/it is sufficient to show) for penultimate steps
		\end{itemize}
	\end{strategybox}

	\subsection{Additional Strategic Approaches}
	
	\subsubsection{Draw a Picture! (Section 2.4)}

\begin{strategybox}{Draw a Picture!}
	\textbf{Visual Translation Strategy:}
	\begin{itemize}
		\item Convert algebraic or numerical problems into geometric representations
		\item Use distance-time graphs, coordinate systems, or vector diagrams
		\item Interpret variables as coordinates or geometric objects
	\end{itemize}
	
	\textbf{Key Applications:}
	\begin{itemize}
		\item Distance-time graphs reveal intersection points naturally
		\item Lattice point problems become counting exercises
		\item Vector interpretation simplifies constraint problems
		\item "A picture is worth a thousand calculations"
	\end{itemize}
	
	\textbf{Examples:}
	\begin{itemize}
		\item Monk problem: Distance-time graph shows paths must intersect
		\item Ordered pairs problem: Count lattice points in parallelogram
		\item Geometric series: Visualize as areas or lengths
	\end{itemize}
\end{strategybox}

	\subsubsection{Recast the Problem (Section 2.4)}

\begin{strategybox}{Recast the Problem}
	\textbf{Cross-Domain Translation:}
	\begin{itemize}
		\item Geometric $\leftrightarrow$ Algebraic (analytic geometry)
		\item Combinatorial $\leftrightarrow$ Number theoretic
		\item Logical $\leftrightarrow$ Geometric (coloring arguments)
		\item Numerical $\leftrightarrow$ Symbolic
	\end{itemize}
	
	\textbf{Coloring as Reformulation:}
	\begin{itemize}
		\item Introduce colors to reveal invariants or constraints
		\item Chessboard coloring: Domino tiling problem becomes parity argument
		\item Gallery problem: Graph coloring yields elegant solution
		\item Convert spatial problems to counting problems
	\end{itemize}
	
	\textbf{Strategic Principle:}
	\begin{itemize}
		\item "Open your mind to other ways of reinterpreting problems"
		\item Implementation is tactical, requiring specialized knowledge
		\item See Chapter 4 for detailed crossover tactics
	\end{itemize}
\end{strategybox}

	\subsubsection{Change Your Point of View (Section 2.4)}

\begin{strategybox}{Change Your Point of View}
	\textbf{Perspective Shifting:}
	\begin{itemize}
		\item Choose the "natural" point of view for the problem
		\item Consider perspectives of different actors/objects in the problem
		\item Look for reference frames that simplify analysis
		\item Manifestation of mathematical peripheral vision
	\end{itemize}
	
	\textbf{Classic Example - Swimming Problem:}
	\begin{itemize}
		\item Standard view: complex relative motion calculations
		\item Hat's point of view: swimmer travels equal time in both directions
		\item Result: 2-hour round trip, 1 mile traveled $\Rightarrow$ current is 0.5 mph
	\end{itemize}
	
	\textbf{Related Strategies:}
	\begin{itemize}
		\item Combines with symmetry (Four Bugs problem, Example 3.1.6)
		\item Often reveals hidden simplicity in apparently complex problems
		\item Particularly powerful in relative motion and optimization problems
	\end{itemize}
   \end{strategybox}
	
	\section{Tactical Methodology Breakdown}
	
	\subsection{Core Mathematical Tactics}
	
	\subsubsection{Symmetry (Chapter 3.1)}
	
	\begin{tacticsbox}{Symmetry}
		\textbf{Geometric Symmetry:}
		\begin{itemize}
			\item Rotational and reflectional symmetries
			\item Coordinate transformation techniques
			\item Symmetry-preserving operations
		\end{itemize}
		
		\textbf{Algebraic Symmetry:}
		\begin{itemize}
			\item Symmetric polynomials and expressions
			\item Variable substitution using symmetry
			\item WLOG arguments based on symmetric cases
		\end{itemize}
	\end{tacticsbox}
	
	\subsubsection{The Extreme Principle (Chapter 3.2)}
	
	\begin{tacticsbox}{The Extreme Principle}
		\textbf{Methodology:}
		\begin{itemize}
			\item Consider extreme cases (maximum/minimum values)
			\item "Maximize balanced wires" example from Affirmative Action Problem
			\item Applications in optimization and existence proofs
		\end{itemize}
		
		\textbf{Implementation Strategy:}
		\begin{enumerate}
			\item Identify quantity to extremize
			\item Assume extremal configuration exists
			\item Derive properties of extremal case
			\item Use contradiction if extremal case leads to impossibility
		\end{enumerate}
	\end{tacticsbox}
	
	\subsubsection{The Pigeonhole Principle (Chapter 3.3)}
	
	\begin{tacticsbox}{The Pigeonhole Principle}
		\textbf{Three Levels of Application:}
		\begin{enumerate}
			\item \textbf{Basic:} $n+1$ objects in $n$ boxes $\Rightarrow$ some box has $\geq 2$ objects
			\item \textbf{Intermediate:} Generalized forms with unequal distributions
			\item \textbf{Advanced:} Abstract applications to infinite sets and continuous domains
		\end{enumerate}
		
		\textbf{Problem-Solving Approach:}
		\begin{itemize}
			\item Identify "pigeons" (objects to be distributed)
			\item Identify "pigeonholes" (categories or boxes)
			\item Count carefully to establish the pigeonhole condition
		\end{itemize}
	\end{tacticsbox}
	
	\subsubsection{Invariants (Chapter 3.4)}
	
	\begin{tacticsbox}{Invariants}
		\textbf{Types of Invariants:}
		\begin{enumerate}
			\item \textbf{Parity:} Even/odd properties preserved under operations
			\item \textbf{Modular Arithmetic:} Remainder classes modulo $n$
			\item \textbf{Monovariants:} Quantities that strictly increase/decrease
		\end{enumerate}
		
		\textbf{Strategic Application:}
		\begin{itemize}
			\item Identify operations that preserve certain properties
			\item Use invariants to prove impossibility results
			\item Apply coloring arguments based on invariant preservation
		\end{itemize}
	\end{tacticsbox}

	\subsection{Crossover Tactics}
	
	\subsubsection{Graph Theory Applications (Chapter 4.1)}
	
	\begin{tacticsbox}{Graph Theory Applications}
		\textbf{Key Concepts:}
		\begin{itemize}
			\item Connectivity and cycle detection
			\item Eulerian and Hamiltonian paths
			\item Graph coloring and planar graphs
		\end{itemize}
		
		\textbf{Problem Translation:}
		\begin{itemize}
			\item Convert relationship problems to graph problems
			\item Use vertex/edge properties to model constraints
			\item Apply graph algorithms to solve combinatorial problems
		\end{itemize}
	\end{tacticsbox}
	
	\subsubsection{Complex Numbers (Chapter 4.2)}
	
	\begin{tacticsbox}{Complex Numbers}
		\textbf{Geometric Applications:}
		\begin{itemize}
			\item Rotation and scaling transformations
			\item Roots of unity in regular polygons
			\item Complex plane as powerful geometric tool
		\end{itemize}
		
		\textbf{Algebraic Techniques:}
		\begin{itemize}
			\item Polynomial root analysis
			\item Trigonometric identity derivation
			\item Geometric series and complex exponentials
		\end{itemize}
	\end{tacticsbox}
	
	\subsubsection{Generating Functions (Chapter 4.3)}
	
	\begin{tacticsbox}{Generating Functions}
		\textbf{Applications:}
		\begin{itemize}
			\item Recurrence relation solutions
			\item Combinatorial counting problems
			\item Partition theory and number theory
		\end{itemize}
		
		\textbf{Methodology:}
		\begin{enumerate}
			\item Define the problem sequence or structure to be analyzed
			\item Encode sequence as power series
			\item Manipulate generating function algebraically
			\item Extract coefficients to solve original problem
		\end{enumerate}
	\end{tacticsbox}
	
	\section{Problem Classification System}
	
	\subsection{Problem Types by Structure}
	
	\subsubsection{Problem Categories}
	
	\begin{examplebox}{Problem Categories}
		\textbf{By Solution Type:}
		\begin{enumerate}
			\item \textbf{"To Find" Problems:} Seek specific numerical or algebraic answers
			\item \textbf{"To Prove" Problems:} Require general logical arguments
		\end{enumerate}
		
		\textbf{By Source and Difficulty:}
		\begin{enumerate}
			\item \textbf{Recreational Problems:} Brain teasers requiring minimal formal mathematics
			\item \textbf{Contest Problems:} Formal examination problems with time constraints
			\begin{itemize}
				\item AHSME/AIME level (high school)
				\item USAMO/IMO level (olympiad)
				\item Putnam level (undergraduate)
			\end{itemize}
			\item \textbf{Open-Ended Problems:} Possibly unsolvable investigation problems
		\end{enumerate}
	\end{examplebox}
	
	\subsection{Difficulty Progression Framework}
	
	\subsubsection{Pedagogical Progression}
	
	\begin{examplebox}{Pedagogical Progression}
		\textbf{Skill Development Sequence:}
		\begin{enumerate}
			\item \textbf{Orientation Skills:} Problem reading and classification
			\item \textbf{Investigation Techniques:} Experimental and exploratory methods
			\item \textbf{Pattern Recognition:} Conjecture formation and testing
			\item \textbf{Proof Construction:} Rigorous argument development
			\item \textbf{Generalization:} Extension to broader contexts
		\end{enumerate}
		
		\textbf{Mathematical Maturity Indicators:}
		\begin{itemize}
			\item Comfort with multiple solution approaches
			\item Ability to recognize when problems are related
			\item Skill in adapting known techniques to novel situations
			\item Persistence through extended investigations
		\end{itemize}
	\end{examplebox}
	
	\section{Interconnections and Meta-Framework}
	
	\subsection{Strategy-Tactics-Tools Relationships}
	
	\subsubsection{Integration Principles}
	
	The book emphasizes that successful problem solving requires fluid movement between all three levels:
	
	\begin{strategybox}{Integration Principles}
		\textbf{Hierarchical Flow:}
		\begin{itemize}
			\item Strategies guide overall approach and investigation direction
			\item Tactics provide concrete methods for implementing strategic plans
			\item Tools execute specific computational or logical steps
		\end{itemize}
		
		\textbf{Bidirectional Feedback:}
		\begin{itemize}
			\item Tool-level failures may require tactical adjustments
			\item Tactical difficulties may necessitate strategic reorientation
			\item Success at any level informs approach at other levels
		\end{itemize}
	\end{strategybox}
	
	\subsection{Cross-Domain Applications}
	
	\subsubsection{Mathematical Unity}
	
	\begin{tacticsbox}{Mathematical Unity}
		\textbf{Recurring Themes:}
		\begin{itemize}
			\item Symmetry appears in algebra, geometry, and combinatorics
			\item Extremal principles apply across all mathematical domains
			\item Invariant techniques work in discrete and continuous settings
		\end{itemize}
		
		\textbf{Transfer Skills:}
		\begin{itemize}
			\item Graph theory models geometric and algebraic relationships
			\item Complex numbers bridge algebra and geometry
			\item Generating functions connect combinatorics and analysis
		\end{itemize}
	\end{tacticsbox}
	
	\section{Educational Implementation Guidelines}
	
	\subsection{Instructional Recommendations}
	
	\subsubsection{Teaching Methodology}
	
	\begin{examplebox}{Teaching Methodology}
		\textbf{Active Learning Emphasis:}
		\begin{itemize}
			\item "Read with pencil and paper by your side" (p. 10)
			\item Attempt each example before reading solutions
			\item Maintain multiple unsolved "backburner" problems
		\end{itemize}
		
		\textbf{Progressive Skill Building:}
		\begin{enumerate}
			\item Start with orientation and investigation strategies
			\item Gradually introduce tactical methods
			\item Build tool proficiency through practice
			\item Develop integration skills through complex problems
		\end{enumerate}
	\end{examplebox}
	
	\subsection{Assessment and Evaluation}
	
	\subsubsection{Evaluation Framework}
	
	\begin{strategybox}{Evaluation Framework}
		\textbf{Process-Oriented Assessment:}
		\begin{itemize}
			\item Value investigation effort over solution completeness
			\item Assess strategic thinking and approach selection
			\item Evaluate communication and reasoning clarity
		\end{itemize}
		
		\textbf{Growth Indicators:}
		\begin{itemize}
			\item Increased persistence on challenging problems
			\item Improved pattern recognition across domains
			\item Enhanced ability to select appropriate techniques
			\item Greater confidence in exploring novel approaches
		\end{itemize}
	\end{strategybox}
	
	\section{Conclusion}
	
	Paul Zeitz's framework represents a comprehensive approach to mathematical problem solving that balances rigor with creativity, structure with flexibility. The three-tier hierarchy provides clear organizational principles while maintaining the exploratory spirit essential to mathematical discovery.
	
	\begin{strategybox}{Key Takeaways}
		\textbf{Fundamental Principles:}
		\begin{enumerate}
			\item Problem solving is learnable through systematic practice
			\item Psychological factors are as important as mathematical knowledge
			\item Investigation and exploration precede formal argumentation
			\item Mathematical domains are unified by common strategic approaches
		\end{enumerate}
		
		\textbf{Implementation Success Factors:}
		\begin{itemize}
			\item Emphasis on process over product
			\item Integration of multiple mathematical domains
			\item Progressive skill development with appropriate scaffolding
			\item Cultivation of mathematical confidence and persistence
		\end{itemize}
	\end{strategybox}
	
	The framework serves as both pedagogical guide and practical problem-solving resource, suitable for implementation at high school through undergraduate levels with appropriate modifications for student mathematical maturity and instructional context.
	
\end{document}
